\documentclass[11pt]{article}
\usepackage[utf8]{inputenc}
\usepackage{geometry}
\usepackage{xcolor}
\usepackage{enumitem}

\geometry{
    a4paper,
    margin=1in
}

% Define the pink accent color used in the heading
\definecolor{creativepink}{HTML}{E6007E}

\begin{document}

\begin{flushleft}
    {\huge\bfseries\color{creativepink}\itshape 6. Creative Writing}
\end{flushleft}

\vspace{1.5em}

\noindent \textsc{The} teacher was explaining the lines in the beginning of Shakespeare’s play \textit{Macbeth}. It was a description of the battle and the lines were:

\vspace{1em}
\begin{quote}
    Like Valour’s minion, carved out his passage,\\
    Till he faced the slave;\\
    With ne’er shook hands, nor baded farewell to him,\\
    Till he \textbf{unseam’d} him from the nave to the chaps,...
\end{quote}
\vspace{1em}

\noindent The teacher asked the students what the word ‘unseamed’ meant. It was difficult. The teacher prodded them on. “What does ‘seam’ mean? Haven’t you ever come across the word?” One of the students blurted out “Cricket ball”.

\vspace{1em}
\noindent This is an example of how each of us reacts to words according to what our own experience has been.

\vspace{1em}
\noindent When we write about factual information, all of us write almost similarly. But when we write for pleasure each of us may write about the same event in different ways.

\vspace{1em}
\noindent \textit{One very important element in creative writing is imagination. This is reflected in}

\begin{itemize}[itemsep=0.2em, topsep=0.5em]
    \item our view or perspective
    \item choice of words
    \item the comparisons we make
    \item the images we use
    \item the tone we adopt
    \item novelty of ideas.
\end{itemize}

\vspace{1em}
\noindent Let us study the paragraph below.

\vspace{1em}
\noindent A town is like an animal. A town has a nervous system and a head and shoulders and feet. A town is a thing separate from all other towns, so that there are no towns alike. And a town has a whole emotion. How news travels through a town

\vfill
\begin{center}
    \small 2024-25
\end{center}

\end{document}
\documentclass[11pt]{article}
\usepackage[utf8]{inputenc}
\usepackage{geometry}
\usepackage{fancyhdr}

\geometry{
    a4paper,
    margin=1in
}

% Header configuration matching the textbook layout
\pagestyle{fancy}
\fancyhf{}
\fancyhead[L]{\small\scshape Creative Writing}
\fancyhead[R]{\small 103}
\renewcommand{\headrulewidth}{0pt}
\setlength{\headheight}{14pt}

\begin{document}

\noindent is a mystery not easily to be solved. News seems to move faster than small boys can scramble and dart to tell it, faster than women can call it over the fences. (from an adapted version of Steinbeck’s \textit{The Pearl})

\vspace{1em}
\noindent \textit{The topic}: A Town \\
\textit{Analogy or comparison}: to an animal \\
\textit{Word choice}: “has a whole emotion.” \\
\textit{Comparisons}: “faster than small boys can scramble and dart, faster than women....”

\vspace{1em}
\noindent We find the first element of imagination operating in the way the writer visualises the town. Then he extends the primary analogy. The tone he adopts is light humour, a little sarcastic.

\vspace{1em}
\noindent When we begin to write a story or poem we let our imagination free. We try to say things in a new way. This novelty is what makes our writing pleasurable to the reader.

\vspace{1em}
\noindent Sometimes sentence structures are also different from factual writing. Consider the following:
\begin{quote}
    They waited in their chairs until the pearls came in, \textbf{and then they cackled and fought and shouted and threatened} until they reached the lowest price the fisherman would stand. (from \textit{The Pearl}).
\end{quote}

\noindent In a normal construction we will not use so many ‘ands’. But the action of the story is best reflected through this kind of chaining of actions through ‘ands’. It is appropriate to the movement of the action described.

\vspace{1.5em}
\noindent Let us look at another example:
\begin{quote}
    She dragged me after her into Miss Rachel’s sitting-room, which opened to her bedroom. At her bedroom door stood Miss Rachel, \textbf{her face almost white as the white dressing-gown she wore.}
\end{quote}

\noindent The author has used a simile: “white as the white dressing-gown she wore.”

\vspace{1em}
\noindent In fact, the whiteness of a human face is because of a strong emotion — fear or shock.

\vspace{1em}
\noindent But here comparing the whiteness to the dressing-gown she wore serves to \textbf{exaggerate} and intensify the emotion.

\end{document}
\documentclass[11pt]{article}
\usepackage[utf8]{inputenc}
\usepackage{geometry}
\usepackage{fancyhdr}
\usepackage{array}

\geometry{
    a4paper,
    margin=1in
}

% Header configuration matching the textbook layout
\pagestyle{fancy}
\fancyhf{}
\fancyhead[L]{\small 104}
\fancyhead[R]{\small\scshape Hornbill}
\renewcommand{\headrulewidth}{0pt}
\setlength{\headheight}{14pt}

\begin{document}

\noindent Exaggeration is one of the ways in which fact is distinguished from fiction.

\vspace{1em}
\noindent Now look at these lines from a well-known poem, ‘An Elegy Written in a Country Churchyard’ by Thomas Gray.

\vspace{1em}
\begin{quote}
    Full many a gem of purest ray serene\\
    The dark unfathom’d caves of ocean bear\\
    Full many a flower is born to blush unseen\\
    And waste its fragrance in the desert air.
\end{quote}
\vspace{1em}

\noindent The stanza carries a simple statement: many people with outstanding qualities live and die unnoticed by the world.

\vspace{1em}
\noindent To state this, the poet has used two strong images, ‘a gem’ and ‘a flower’.

\vspace{1em}
\noindent He has used two contrasting places: the ocean, that is full of water and the desert with no water at all.

\vspace{1em}
\noindent Also notice the rhyming words: ‘serene’ and ‘unseen’, ‘bear’ and ‘air’.

\vspace{1em}
\noindent The first and third lines also begin with the same words — “Full many a”. The lines are of equal length.

\vspace{1em}
\noindent All this together contribute to the literary quality of these lines.

\vspace{2em}
\noindent \textbf{Activity I}

\vspace{0.5em}
\noindent Put down the images that come to your mind immediately when you see the words in the box.

\vspace{1em}
\begin{center}
    \begin{tabular}{|ccccc|}
        \hline
        cat & cupboard & wall & pond & bird \\
        \hline
    \end{tabular}
\end{center}
\vspace{1.5em}

\noindent \textbf{Activity II}

\vspace{0.5em}
\noindent Try to write four lines of poetry or four sentences of prose with one of these as the starting point.

\vspace{1.5em}
\noindent \textbf{Activity III}

\vspace{0.5em}
\noindent Write a short story beginning with this sentence:
\begin{quote}
    When the last of the guests left, I went back into the hall\dots
\end{quote}

\vspace{1.5em}
\noindent \textbf{Activity IV}

\vspace{0.5em}
\noindent Look for a story, a poem and a newspaper article on environment conservation and see how the style of each is different from the other.

\end{document}

